\section{Discussion and Future Scope}
\label{sec:discussion}

The results answer the questions that close Section~\ref{sec:lit}: the published accuracy owed much
to a leak, transfer bought a fast start but not a better end point, and the data show imagined
movement rather than a sixth-finger command. This section reads each answer, states what limits
them and says what should come next.

\subsection{A fast start from one transfer configuration}
One trial of each movement (\rqTwoMinutesKOne\,min) was enough for re-centred transfer to match
the three-class accuracy that the same classifier without transfer reached only after ten
(\rqTwoMinutesKTen\,min nominal, \rqTwoMinutesRecordedKTen\,min as recorded). Beyond that the
transfer curve barely rose, and with the whole pool the patient-only decoder finished numerically
ahead, by a margin we left untested. Earlier work found that transfer cuts calibration for
healthy users~\cite{zanini2018,he2020,wu2022review} and, across sessions and patients, after
stroke~\cite{giles2022,cao2021}; here, transfer between patients paid off during the first few
trials. These curves belong to one configuration and are no verdict on transfer as such: the
target weight let a single window of the patient count no more than one of the thousands from
other patients, it won a selection in which every candidate scored close to chance, and the deep
pair, too, ended with its least adaptive settings. The ablations add that re-centring made the early lead on sixth finger against rest, while on
three classes pooling alone gave about half of it, and part of the lead reflects how little any
decoder learns from four windows. A weight that grows with $k$, a hand-over to the patient's own decoder when
cross-validation inside the calibration data prefers it, or Procrustes
alignment~\cite{rodrigues2019} might keep transfer improving; we tried none of them.

\subsection{Leakage}
The inflated figure comes from the order of the steps, not from the data: the same code inflated
accuracy on two healthy datasets, and the check was possible only because the authors released
both data and code. Our rebuild did not reach the published figure, yet each of the three
leakage-free estimates sat well below the leaky one. Anyone reusing these recordings should keep
every supervised step inside the training folds, separate the data by run, and quote the
leakage-free number next to the published one. Measured that way, the two imagined movements
differ only weakly in this EEG, and some of the difference may be visual.

\subsection{Imagery against rest, not yet a sixth finger}
The study neither showed nor excluded a separate sixth-finger command. In the planned control, a
decoder never shown a sixth-finger trial recognised sixth-finger imagery at least as well, and
every detector also answered imagined hand movements. The mu rhythm at C3 fell a little further
for the sixth finger (Section~\ref{sec:method}), but this was seen before the freeze and the plan
ruled out testing it again; a difference in depth would not by
itself make the movements distinguishable. Offline, the data support one contrast, imagery against
rest, which could serve as a ``move now'' switch. Replayed continuously it fired in under half of
the sixth-finger periods, late in each period and for the hand too, so only a closed-loop study
can say whether a patient could operate it. It would need recalibrating each day; the earlier
session plus five new trials per movement is how the fast start would look in practice.

\subsection{Limitations}
All \nSubjects{} patients were treated at one centre, and between sessions they received one of
two different therapies. Every evaluation was offline or replayed; no patient used the decoder in
closed loop, and the sixth finger was worn during imagery but never moved by a decoder. The
patient-only baselines are handicapped at small budgets by design, with four windows at $k = 1$
and, for the new command, unbalanced classes without class weights; a regularised or
class-weighted baseline could shrink transfer's early lead. The plan itself declared two limits:
rest windows can contain the beta rebound that follows imagery, which makes any accuracy involving
rest an upper bound for self-paced rest, and the sixth finger's extra desynchronisation is spread
widely, so it may track how intensely the patient imagined rather than a separate motor map. Rest
always came with a cross on screen and imagery with a picture; the switch relies mainly on central
channels, yet a visual share cannot be excluded, and the hand-against-finger information is at
least as strong at the back of the head. Some of the authors' cleaning was done by hand and cannot
be repeated exactly, and aligning the published windows with the raw files left a few windows and
one patient's run order unresolved. The plan was frozen in git rather than with an independent
registry, and the code is shared on request, not openly. Finally, \nPre{} patients give power only
for medium or large effects; the inconclusive H3 and H3c do not show that no effect exists.

\subsection{Future scope}
The decisive next step is a closed-loop study in which the decoder moves the finger. For
calibration, the untested options above and a class-weighted or shrinkage baseline would show how
much of the plateau and the early lead is due to the configuration. A check resolved in time,
comparing the first second after the cue with the later ones, and a paradigm that keeps the same
screen for rest and imagery would tackle the visual confound. Further stroke
datasets~\cite{liu2024acute,lu2026chronic} would test the results beyond one hospital, and EEG
foundation models~\cite{jiang2024labram,wang2025cbramod} might lift the plateau, which
pre-training on healthy grasp imagery failed to do in a first try. Registering the plan
independently and publishing the code would let others verify this work.
